\documentclass{article}

%---------------------------------------------------------------------------
%  Packages needed by nearly every document
%---------------------------------------------------------------------------

% No need to worry about possibility of packages being loaded multiple
% times.  LaTeX has built-in ability to detect this and load only once

% Title and section formatting
\usepackage[small,sf,bf]{titlesec}
%\titleformat*{\section}{\Large\bfseries\sffamily}
%\titleformat*{\subsection}{\large\bfseries\sffamily}

% Math tools
\usepackage{amsmath,amssymb,amsthm,mathtools}
\usepackage{bigints}

% Better access of bold math symbols
\usepackage{bm}
% Allows the striking of terms in algebraic expressions
\usepackage{cancel}

\DeclarePairedDelimiter\set\{\}

% Libertine fonts
\usepackage[ttscale=0.75]{libertine}
\usepackage[libertine]{newtxmath}

% Including figures
\usepackage{graphicx}
\graphicspath{{Figures}}

% Enhanced formatting of dates and times
\usepackage{datetime2}
\renewcommand{\DTMdisplaydate}[4]{#1-#2-#3}

% Automatic adjustment of table column widths to achieve
% an overall table width
\usepackage{tabularx}
% Better formatting of tables (\toprule, \midrule and \bottomrule)
\usepackage{booktabs}

% Sideways oriented tables and figures
\usepackage{rotating}

% TeX commands for drawing graphics
\usepackage{tikz}

% Enhanced handling of figure and table captions
\usepackage{caption}

% Specify the counter format easily for enumerate environments
\usepackage{enumerate}

% Enables bold math Greek symbols
\usepackage{bm}

% Next lines set up the system for formatting units
\usepackage{siunitx}
\sisetup{detect-weight=true, detect-family=true}
\sisetup{range-phrase={ to },%
	range-units =brackets,%
	table-number-alignment =center,%
	table-figures-exponent=1%
}

\DeclareSIUnit\atm{atm}
\DeclareSIUnit\bar{bar}
\DeclareSIUnit\atom{atom}

% Next lines setup hyperlink formatting
\usepackage[colorlinks=true,
	linkcolor=blue,
	urlcolor=blue,
	breaklinks,
	pdftex,
	allbordercolors=white]{hyperref}

% Can convert all text in environment to upper case or lower case
\usepackage{textcase}

% Chemical symbols, formulas, and reactions
\usepackage{chemformula}

%Fancy lists
\usepackage{listings}

% Enhanced color commands
\usepackage{xcolor}

% Definition of some useful colors
\definecolor[named]{darkblue}{rgb}{0,0,0.7}
\definecolor[named]{darkgreen}{rgb}{0,0.5,0}
\definecolor[named]{aggiemaroon}{rgb}{0.3137,0,0}
\definecolor[named]{aggiegray}{rgb}{0.2,0.1725,0.1725}

% Format for program blocks
%\lstset{frame=tb,
%  language=C,
%  aboveskip=3mm,
%  belowskip=3mm,
%  showstringspaces=false,
%  columns=flexible,
%  basicstyle={\small\ttfamily},
%  numbers=none,
%  numberstyle=\tiny\color{gray},
%  keywordstyle=\color{blue},
%  commentstyle=\color{dkgreen},
%  stringstyle=\color{mauve},
%  breaklines=true,
%  breakatwhitespace=true,
%  tabsize=3
%}

% Small vertical space between paragraphs, no indentation of first line
\usepackage{parskip}

% Number equations by section
\numberwithin{equation}{subsection}

\usepackage{etoolbox}
\makeatletter
\patchcmd{\@startsection}% <cmd>
{\if@noskipsec}% <search>
{\csgdef{theequation}{\csname the#1\endcsname.\arabic{equation}}\if@noskipsec}% <replace>
{}{}% <success><failure>
\makeatother

\usepackage[backend=biber,
	sorting=none,%
	sortcites,%
	maxnames=4,%
	citestyle=authoryear,%
	uniquename=init,%
	giveninits=true,%
	bibstyle=authoryear,%
	terseinits=true,%
	casechanger=expl3,%
]{biblatex}
\addbibresource{/Users/bullard/Library/texmf/bibtex/bib/master.bib}

\DeclareFieldFormat
[article,inbook,incollection,inproceedings,patent,thesis,unpublished]
{titlecase}{\MakeSentenceCase*{#1}}


% The package xparse helps handle complex optional arguments to new commands
\usepackage{xparse}
\usepackage{ifmtarg}
\usepackage{textcase}
\usepackage{bm}
\usepackage{chemformula}

% These commands use the xparse package

% Symbols for useful constants
\NewDocumentCommand{\gasconst}{s}{%
	\IfBooleanTF{#1}%
	{\ensuremath{\num{8.314}}} {\ensuremath{R}}%
}
\NewDocumentCommand{\boltz}{s}{%
	\IfBooleanTF{#1}%
	{\ensuremath{\num{1.38e-23}}} {\ensuremath{k_B}}%
}
\NewDocumentCommand{\planck}{s}{%
	\IfBooleanTF{#1}%
	{\ensuremath{\num{6.63e-34}}} {\ensuremath{h}}%
}
\NewDocumentCommand{\avogadro}{s}{%
	\IfBooleanTF{#1}%
	{\ensuremath{\num{6.022e23}}} {\ensuremath{N_{A}}}%
}

% Declaration of non-standard units
\DeclareSIUnit\bar{bar}
\DeclareSIUnit\year{y}

\NewDocumentCommand{\RT}{}{\ensuremath{\gasconst T}}
\NewDocumentCommand{\kT}{}{\ensuremath{\boltz T}}

% Common decorations
\NewDocumentCommand{\cmp}{}{\ensuremath{*}}
\NewDocumentCommand{\cmplx}{}{\ensuremath{*}}
\NewDocumentCommand{\std}{}{\ensuremath{\circ}}
\NewDocumentCommand{\pmolar}{m}{\ensuremath{\overline{#1}}}

% Thermodynamic potentials
% Upper case is molar, lower case is volumetric

\NewDocumentCommand{\Gibbs}{s}{%
	\IfBooleanTF{#1}%
	{\ensuremath{\mathbb{G}}} {\ensuremath{\text{G}}}%
}
\NewDocumentCommand{\gibbs}{}{\ensuremath{\text{g}}}

\NewDocumentCommand{\Helmholtz}{s}{%
	\IfBooleanTF{#1}%
	{\ensuremath{\mathbb{A}}} {\ensuremath{\text{A}}}%
}
\NewDocumentCommand{\helmholtz}{}{\ensuremath{\text{a}}}

\NewDocumentCommand{\Entropy}{s}{%
	\IfBooleanTF{#1}%
	{\ensuremath{\mathbb{S}}} {\ensuremath{\text{S}}}%
}
\NewDocumentCommand{\entropy}{}{\ensuremath{\text{s}}}

\NewDocumentCommand{\Energy}{s}{%
	\IfBooleanTF{#1}%
	{\ensuremath{\mathbb{U}}} {\ensuremath{\text{U}}}%
}
\NewDocumentCommand{\energy}{}{\ensuremath{\text{u}}}

\NewDocumentCommand{\Work}{s}{%
	\IfBooleanTF{#1}%
	{\ensuremath{\mathbb{W}}} {\ensuremath{\text{W}}}%
}
\NewDocumentCommand{\Enthalpy}{s}{%
	\IfBooleanTF{#1}%
	{\ensuremath{\mathbb{H}}} {\ensuremath{\text{H}}}%
}

\NewDocumentCommand{\Vol}{s}{%
	\IfBooleanTF{#1}%
	{\ensuremath{\mathbb{V}}} {\ensuremath{V}}%
}
\NewDocumentCommand{\Volume}{s}{%
	\IfBooleanTF{#1}%
	{\ensuremath{\mathbb{V}}} {\ensuremath{V}}%
}

\NewDocumentCommand{\area}{s}{%
	\IfBooleanTF{#1}%
	{\ensuremath{\mathbb{A}}} {\ensuremath{A}}%
}

\NewDocumentCommand{\molecvol}{}{\ensuremath{\Omega}}

\NewDocumentCommand{\enthalpy}{}{\ensuremath{\text{h}}}
\NewDocumentCommand{\heat}{}{\ensuremath{q}}
\NewDocumentCommand{\Heat}{}{\ensuremath{Q}}

\NewDocumentCommand{\chempot}{}{\ensuremath{\mu}}
\NewDocumentCommand{\Press}{}{\ensuremath{P}}
\NewDocumentCommand{\ppress}{}{\ensuremath{p}}

\NewDocumentCommand{\Eact}{}{\ensuremath{E_{a}}}
\NewDocumentCommand{\acoeff}{}{\ensuremath{\gamma}}
\NewDocumentCommand{\actprod}{}{\ensuremath{Q}}
\NewDocumentCommand{\eqconst}{}{\ensuremath{K}}
\NewDocumentCommand{\fug}{}{\ensuremath{f}}
\NewDocumentCommand{\molfrac}{}{\ensuremath{X}}
\NewDocumentCommand{\moles}{}{\ensuremath{n}}
\NewDocumentCommand{\molecules}{}{\ensuremath{N}}
\NewDocumentCommand{\nsites}{}{\ensuremath{N}}
\NewDocumentCommand{\nvac}{}{\ensuremath{N_{\text{v}}}}

% Concentration
% If optional star token is given, puts brackets around optional string
% If optional star token is missing, ignores optional string and just types c
% Subscripts and superscripts are handled with the dec command below
\NewDocumentCommand{\conc}{sO{\cdots}}{%
	\IfBooleanTF{#1}%
	{\ensuremath{\left[#2\right]}} {\ensuremath{c}}%
}

% Activity
% If optional star token is given, puts curly braces around optional string
% If optional star token is missing, ignores optional string and just types a
% Subscripts and superscripts are handled with the dec command below
\NewDocumentCommand{\act}{sO{\cdots}}{%
	\IfBooleanTF{#1}%
	{\ensuremath{\set{#2}}} {\ensuremath{a}}%
}

\NewDocumentCommand{\oderivh}{O{}mm}{\ensuremath{\text{d}^{#1}#2/\text{d}{#3}^{#1}}}
\NewDocumentCommand{\oderivv}{O{}mm}{\ensuremath{\frac{\text{d}^{#1}#2}{\text{d}{#3}^{#1}}}}
\NewDocumentCommand{\pderivv}{O{}mm}{\ensuremath{\frac{\partial^{#1} #2}{\partial #3^{#1}}}}
\NewDocumentCommand{\pderivh}{O{}mm}{\ensuremath{\partial^{#1} #2/\partial #3^{#1}}}
\NewDocumentCommand{\rateconst}{}{\ensuremath{k}}
\NewDocumentCommand{\Energ}{}{\ensuremath{E}}
\NewDocumentCommand{\energ}{}{\ensuremath{\epsilon}}
\NewDocumentCommand{\DOS}{}{\ensuremath{\omega}}
\NewDocumentCommand{\exval}{m}{\ensuremath{\langle #1 \rangle}}
\NewDocumentCommand{\volfrac}{}{\ensuremath{\phi}}
\NewDocumentCommand{\mass}{}{\ensuremath{m}}
\NewDocumentCommand{\molarmass}{}{\ensuremath{M}}
\NewDocumentCommand{\flux}{}{\ensuremath{J}}
\NewDocumentCommand{\fluxp}{}{\ensuremath{\flux^{\prime}}}
\NewDocumentCommand{\diffus}{}{\ensuremath{D}}
\NewDocumentCommand{\dens}{}{\ensuremath{\rho}}

% Symbol for nucleation rate
\NewDocumentCommand{\nucrate}{}{\ensuremath{\Theta}}

% Notation for vectors
% The * token will make a unit vector
% The + token will indicate the magnitude of the vector
% Both tokens can be used together, but * must come before +

\NewDocumentCommand{\vect}{st+m}{%
	\IfBooleanTF{#1} {%
		\IfBooleanTF{#2} {%
			\ensuremath{\lvert \mathbf{\hat{#3}} \rvert}%
		}{%
			\ensuremath{\mathbf{\hat{#3}}}%
		}%
	}{%
		\IfBooleanTF{#2} {%
			\ensuremath{\lvert \mathbf{#3} \rvert}%
		}{%
			\ensuremath{\mathbf{#3}}%
		}%
	}%
}

% Shorthand partial derivatives
% The * token produces shorthand notation
%{ \ensuremath{#3_{\myrepeat{#2}{#4}}} }
% Both tokens can be used together, but * must come before +

\ExplSyntaxOn
\NewExpandableDocumentCommand{\myrepeat}{mm}
{
	\int_compare:nT { #1 > 0 }
	{
		#2 \prg_replicate:nn { #1 - 1 } { #2 }
	}
}
\ExplSyntaxOff

\NewDocumentCommand{\pderiv}{somm}{%
	\IfBooleanTF{#1} {%
		\IfNoValueTF{#2}
		{ \ensuremath{#3_{#4}} }
		{ \ensuremath{#3_{\myrepeat{#2}{#4}}} }
	}{%
		\IfNoValueTF{#2}
		{ \ensuremath{\frac{\partial #3}{\partial #4}} }
		{ \ensuremath{\frac{\partial^{#2} #3}{\partial #4^{#2}}} }
	}%
}

\NewDocumentCommand{\oderiv}{somm}{%
	\IfBooleanTF{#1} {%
		\IfNoValueTF{#2}
		{ \ensuremath{#3^{\prime}} }
		{ \ensuremath{#3^{\myrepeat{#2}{\prime}}} }
	}{%
		\IfNoValueTF{#2}
		{ \ensuremath{\frac{\text{d} #3}{\text{d}#4}} }
		{ \ensuremath{\frac{\text{d}^{#2} #3}{\text{d}#4^{#2}}} }
	}%
}

% Shorthand for divergence and gradient operators
\NewDocumentCommand{\diverg}{s}{%
	\IfBooleanTF{#1} {%
		\ensuremath{\text{div}\,}%
	}{%
		\ensuremath{\vect{\nabla} \cdot }%
	}%
}

\NewDocumentCommand{\grad}{s}{%
	\IfBooleanTF{#1} {%
		\ensuremath{\text{grad}\,}%
	}{%
		\ensuremath{\vect{\nabla} }%
	}%
}

\NewDocumentCommand{\curl}{s}{%
	\IfBooleanTF{#1} {%
		\ensuremath{\text{curl}\,}%
	}{%
		\ensuremath{\vect{\nabla} \times }%
	}%
}

\NewDocumentCommand{\mylaplac}{s}{%
	\IfBooleanTF{#1} {%
		\ensuremath{\div \grad\,}%
	}{%
		\ensuremath{\grad^2 }%
	}%
}


\NewDocumentCommand{\surfnrg}{}{\ensuremath{\gamma}}
\NewDocumentCommand{\Grate}{}{\ensuremath{G}}
\NewDocumentCommand{\freq}{}{\ensuremath{\nu}}

% Viscosity
\NewDocumentCommand{\visc}{}{\ensuremath{\mu}}
\NewDocumentCommand{\stress}{}{\ensuremath{\sigma}}
\NewDocumentCommand{\strain}{}{\ensuremath{\epsilon}}
\NewDocumentCommand{\sstress}{}{\ensuremath{\tau}}
\NewDocumentCommand{\sstrain}{}{\ensuremath{\gamma}}
\NewDocumentCommand{\sstrainrate}{}{\ensuremath{\dot{\gamma}}}
\NewDocumentCommand{\shearmod}{}{\ensuremath{G}}
\NewDocumentCommand{\bulkmod}{}{\ensuremath{K}}
\NewDocumentCommand{\freevol}{}{\ensuremath{V}}
\NewDocumentCommand{\tecoeff}{}{\ensuremath{\alpha}}

% Heat capacity
\NewDocumentCommand{\heatcap}{}{\ensuremath{C}}

% Mathematical spaces
\NewDocumentCommand{\Space}{sO{}}{%
	\IfBooleanTF{#1}{%
		\ensuremath{\mathbb{Z}^{#2}}
	}{%
		\ensuremath{\mathbb{R}^{#2}}
	}%
}

% Bold notation for gradient
\NewDocumentCommand{\bgrad}{}{\ensuremath{\pmb{\nabla}}}

% Function notation
\NewDocumentCommand{\fnc}{sO{}m}{%
	\IfBooleanTF{#1}{%
		\ensuremath{\widetilde{#3}_{#2}}
	}{%
		\ensuremath{\tilde{#3}_{#2}}
	}%
}

\theoremstyle{definition}

% Theorem-like environments
\newtheorem{example}{Example}
\newtheorem{definition}{Definition}
\newtheorem{remark}{Remark}

% Simulates filled in text for annotations I might want to keep invisible
% The * token will make the text invisible
% The + token will underline the text
% The default color will be aggiemaroon
% Both tokens can be used together, but * must come before +

\NewDocumentCommand{\fillin}{st+O{aggiemaroon}m}{%
\IfBooleanTF{#1} {%
	\IfBooleanTF{#2}%
	{\underline{\phantom{#4} \hspace{0.25in}}} {\phantom{#4} \hspace{0.25in}}
}{%
	\IfBooleanTF{#2}%
	{\underline{\textcolor{#3}{#4}}} {\textcolor{#3}{#4}}
}%
}

\NewDocumentCommand{\parspace}{}{\vspace{0.125truein}}
\NewDocumentCommand{\dparspace}{}{\vspace{0.25truein}}

% Operators
\NewDocumentCommand{\diff}{}{\ensuremath{\text{d}}}
\NewDocumentCommand{\idiff}{}{\ensuremath{\delta}}

\NewDocumentCommand{\xfact}{O{}}{%
	\IfValueTF{#1}{%
		\ensuremath{\num{#1}\times}%
	}{%
		\ensuremath{\times}%
	}%
}
\NewDocumentCommand{\dotp}{}{\ensuremath{\cdot}}
\NewDocumentCommand{\crossp}{}{\ensuremath{\times}}

% One-stop shop for different decorations for symbols
% Optional argument within < > brackets adds a superscript
% Optional argument within [ ] brackets adds a subscript
% Add a + if you want Delta, or add a * if you want differential
%
% Example 1:  \dec{\Gibbs} --> G
% Example 2:  \dec+{\Gibbs} --> Delta G
% Example 3:  \dec*{\Gibbs} --> d G
% Example 4:  \dec*[i]{\Gibbs} --> d G_i
% Example 5:  \dec*<\std>{\Gibbs} --> d G^o
% Example 5:  \dec+<\std>[i]{\Gibbs} --> Delta G^o_i

\NewDocumentCommand{\dec}{st+d<>d[]m}{%
\IfValueTF{#3}{%
	\IfValueTF{#4}{%
		\IfBooleanTF{#2}
		{\ensuremath{\Delta #5_{#4}^{#3}}} {%
			\IfBooleanTF{#1}
			{\ensuremath{\diff #5_{#4}^{#3}}} {\ensuremath{#5_{#4}^{#3}}}%
		}%
	}{%
		\IfBooleanTF{#2}
		{\ensuremath{\Delta #5^{#3}}} {%
			\IfBooleanTF{#1}
			{\ensuremath{\diff #5_{#4}^{#3}}} {\ensuremath{#5_{#4}^{#3}}}%
		}%
	}%
}{%
	\IfValueTF{#4}{%
		\IfBooleanTF{#2}
		{\ensuremath{\Delta #5_{#4}}} {%
			\IfBooleanTF{#1}
			{\ensuremath{\diff #5_{#4}}} {\ensuremath{#5_{#4}}}%
		}%
	}{%
		\IfBooleanTF{#2}
		{\ensuremath{\Delta #5}} {%
			\IfBooleanTF{#1}
			{\ensuremath{\diff #5}} {\ensuremath{#5}}%
		}%
	}%
}%
}

%----------------------------------------------------------------------------------


\newcommand{\us}{\underline{\phantom{X}}}
\begin{document}

\lstset{basicstyle=\fontfamily{pcr}\scriptsize}

\begin{center}
	\Large{\textbf{\sffamily{THAMES 5.1 User Guide}}}
\end{center}
\begin{center}
	\large{Jeffrey W. Bullard}
\end{center}
\begin{center}
	\large{\DTMnow}
\end{center}

\dparspace
\tableofcontents

\dparspace

\section{\label{sec:requirements} Software Requirements}
The following software must be installed and running properly on your computer:
\begin{itemize}
	\item GEMS3K standalone library 3.4.1: \href{http://gems.web.psi.ch/GEMS3K/downloads/index.html}{http://gems.web.psi.ch/GEMS3K/downloads/index.html}
	\item THAMES v5.1: \href{https://github.tamu.edu/jwbullard/THAMES.git}{https://github.tamu.edu/jwbullard/THAMES.git}
	\item (Optional) GEM-Selektor 3.6.0 or 3.7.0: \href{http://gems.web.psi.ch/GEMS3/downloads/index.html}{http://gems.web.psi.ch/GEMS3/downloads/index.html}
	\item (Optional) ImageMagick 7 or later (for creating image files in PNG format): \href{https://imagemagick.org/}{https://imagemagick.org/}
\end{itemize}

\subsection{\label{sec:path} Putting THAMES in your path}
It will be convenient to add the folder where you install THAMES into your
path. Usually the THAMES binary will be located in the \texttt{bin} subfolder
of your local copy of the git repository that you cloned from GitHub.

\parspace
\begin{example}
	Suppose the THAMES repository is located at \verb|$HOME/thames|. Open
	your \verb|$HOME/.bashrc| or your \verb|$HOME/.zshrc| file, depending
	on whether you use the Bash shell or the Z shell. Somewhere near the bottom
	of that file, add these two lines:
	\small{
		\begin{verbatim}
    PATH="$PATH:$HOME/thames/bin"
    export PATH
  \end{verbatim}
	} \normalsize
	Save and close this file and then run the command
	\small{
		\begin{verbatim}
      source $HOME/.bashrc
  \end{verbatim}
	} \normalsize
	or
	\small{
		\begin{verbatim}
      source $HOME/.zshrc
  \end{verbatim}
	} \normalsize
	You can now test if the system can find the THAMES executable by switching to
	any other directory and run the command \verb|which thames|. This should output
	the location of your THAMES file.
\end{example}

\subsection*{Setting up a sandbox to run THAMES simulations}
It may be convenient to play around with THAMES by using the test cases
provided, but you don't want to pollute your local copy of the repository. Therefore,
it is helpful to make copies of all the test cases in a ``sandbox''. You can
run two commands like these to make a sandbox in your home folder:
\small{
	\begin{verbatim}
    mkdir $HOME/THAMES-Sandbox
    cp -R $HOME/thames/tests/* $HOME/THAMES-Sandbox/.
    \end{verbatim}
} \normalsize
Now let's test whether this setup has worked the way we want by running these
commands:
\small{
	\begin{verbatim}
cd $HOME/THAMES-Sandbox/portcem-298K-sat-wc45
thames
\end{verbatim}
} \normalsize
This command should produce this:
\small{
	\begin{verbatim}
Enter simulation type:
  1) Exit program
  2) Hydration
  3) Leaching
  4) Sulfate attack
\end{verbatim}
} \normalsize
If you see this output, everything is working correctly and you can
\verb|Ctrl-C| to stop.

\section{\label{sec:cmdline} THAMES command line options}
Let's stay in that same \verb|portcem-298K-sat-wc45| folder and run this command:
\small{
	\begin{verbatim}
thames --help 
\end{verbatim}
} \normalsize
This should show the following output:
\small{
	\begin{verbatim}
Usage: "thames [--verbose|-v] [--suppress|-s] [--xyz|-x] [--outfolder|-o] folder
[--help|-h]"
        --verbose [-v]      Produce verbose output
        --suppress [-s]     Suppress warning messages
        --xyz [-x]          Create 3D visualization movie
        --outfolder [-o]    Name of folder for output data (default is Result)

Note: thames --help [-h] will print this help message
\end{verbatim}
} \normalsize

This help message is giving a terse description of the different options you can
include when running THAMES. The options can be added in any order. Each one is
explained a little more fully here.
\begin{itemize}
	\item \verb|--verbose [-v]|: THAMES already produces significant output to your
	      terminal. This output directly to your screen is called the \verb|stdout|
	      stream, which stands for ``standard output''.
	      THAMES writes a lot of information to \verb|stdout| while it runs,
	      including a lot of diagnostic messages (in the future much of
	      this will be removed from normal operation but it is helpful while we
	      are still in development mode.) Issuing this option will create even
	      \textit{more} detailed output. You will usually not need or want this
	      output.
	\item \verb|--suppress [-s]|: This will stop any warning messages from being
	      written to \verb|stdout|. There are not many warning messages, and they can
	      be helpful later to figure out what may have happened if you do not get
	      the results you were expecting. Therefore, it is recommended that you
	      not suppress the warning messages.
	\item \verb|--xyz [-x]|: This option will cause THAMES to create a file with a
	      \verb|.xyz| extension that can be read directly by a 3D visualization
	      application such as Ovito (available for free download at \url{https://ovito.org}).
	      Using this option can cause THAMES to run a little slower because a lot
	      of output is written to this \verb|.xyz| file.
	\item \verb|--outfolder [-o]|: In addition to the stream of output written
	      to \verb|stdout|, THAMES also creates a large number of new output files,
	      and these are all collected in an output folder. This option allows
	      you to specify the name and location of this output folder.
	      \begin{itemize}
		      \item \verb|--outfolder MyFiles| will name the folder \verb|MyFiles| and
		            will put it in the directory where you are running the \verb|thames|
		            command.
		      \item \verb|--outfolder |$\sim$\verb|/SimFiles| will name the folder \verb|SimFiles| but
		            it will be put in your home folder instead of the current folder
		      \item Omitting this option will name the folder \verb|Results| and it will
		            be put in the directory where you are running \verb|thames|.
	      \end{itemize}
\end{itemize}

\subsection{\label{sec:running} Running THAMES non-interactively}
\subsubsection{\label{sec:redirectinginput} Redirecting input}
Running \verb|thames|, with or without any of these options, will enter you into
``interactive'' mode by default, prompting you with menu options on the command
line. When you enter text from your keyboard, you are creating a stream of text
called \verb|stdin|, or the ``standard input'' stream.

In normal operation, THAMES will prompt you for six pieces of information in
this order:
\begin{enumerate}
	\item The type of simulation to run:
	      \small{
		      \begin{verbatim}
      1) Exit program
      2) Hydration
      3) Leaching
      4) Sulfate attack
    \end{verbatim}
	      } \normalsize
	      At the moment options 3 and 4 do not work properly, so you will almost always
	      want to choose ``2''
	\item The name of the master file with thermodynamic information for GEMS.
	\item The name of the file containing the definition of all the microstructure
	      phases that THAMES will recognize for your simulation (JSON format)
	\item The name of the file containing the initial 3D microstructure data
	\item The name of the file containing the list of times at which to calculate
	      a thermodynamic state or to output microstructure data.
	\item The root name that all the output files will have
\end{enumerate}

You can avoid the interactive behavior of THAMES and run a simulation in
non-interactive mode by putting all six of your responses in a text file, each
on a separate line, and then
\textit{redirecting} the input to come from the file instead of fro
\verb|stdin|. Each of the test cases in your sandbox have such a text file
pre-made for you, called \verb|input.in|.

To run THAMES in non-interactive mode, you just need to issue the \verb|thames|
command together with any options, type a space character after the last option,
and then type $<$\verb|input.in &|, and then type \verb|Return| (on Linux or
Windows) or \verb|Enter| on a Mac. So a command that redirects input could be:

\parspace
\small{
	\verb|thames --outfolder MySim| $<$ \verb|input.in| \verb|&|
} \normalsize

Incidentally, the $<$ character in a Unix or Linux command always means ``redirect input
from the file that follows this character''.  Also that last \verb|&| character at the
end of the command is the Unix way of telling the computer to run the command in the
background so you will get your command line prompt back while it is running.

\subsubsection{\label{sec:redirectingoutput} Redirecting output}
If you run the previous command, even in background mode, THAMES will still
write all of its \verb|stdout| stream to your screen, effectively making that
terminal unusable until it finishes. To avoid this behavior, you should always
redirect \verb|stdout| away from your screen and into a text file to be saved
for later. Whereas the symbol for redirecting input is $<$, the symbol for
redirecting all \verb|stdout| is $>$.  However, if something goes wrong and
THAMES issues some kind of error while it is running, that error message will
still go to your terminal. You can direct both normal standard output
(\verb|stdout|) and errors
(called \verb|stderr|) to the same text file by using $>$\verb|&| (with no space
between the two characters) instead of just $>$.

So finally, a full THAMES command that runs in the background non-interactively
would be

\parspace
\small{
	\verb|thames --outfolder MySim| $<$ \verb|input.in| $>$\verb|& output.txt &|
} \normalsize


\section{\label{sec:inputguide} THAMES Input Files}
\subsection{\label{sec:inputcategories} Categories of THAMES Input}
Throughout this document, it will be assumed that your THAMES simulation will happen
in a working directory called \verb!$WORKDIR!.

THAMES requires three categories of input files:
\begin{enumerate}
	\item The initial 3D microstructure of the material system
	\item Thermodynamic data input files produced by GEM-Selektor
	\item Microstructure input file and file linking microstructure phases to GEM phases
	\item Files to specify calculation times and microstructure output times
	\item (Optional) File containing standard input to redirect to the program
\end{enumerate}
Each of these will be detailed in the following sections.

\subsection{\label{sec:microstructureinput} Microstructure initial 3D arrangement}
The microstructure representation is a digitized 3D image with \texttt{X{\us}Size} voxels
in the $x$ direction, \texttt{Y{\us}Size} in the $y$ direction, and
\texttt{Z{\us}Size}
in the $z$ direction.  Each voxel is a physical dimension of
\texttt{Image{\us}Resolution} in micrometer units.
This is a normal ASCII text file.  The first few lines of an example microstructure
are given below.

\scriptsize{
	\begin{lstlisting}
#THAMES:Version: 5.1
#THAMES:X_Size: 100
#THAMES:Y_Size: 100
#THAMES:Z_Size: 100
#THAMES:Image_Resolution: 1.0
1
1
1
2
1
2
1
.
.
.
\end{lstlisting}
}

\normalsize{ }
The meaning of the five lines of header information follows:
\begin{itemize}
	\item \verb!#THAMES:Version! is the THAMES version with which these data
	      should be compatible. It is not currently used in the code but is only for
	      informational purposes;
	\item \verb!#!\texttt{THAMES:X{\us}Size}, \verb!#!\texttt{THAMES:Y{\us}Size}, and \verb!#!\texttt{THAMES:Z{\us}Size}
	      are the system dimensions, in voxel units, in all three Cartesian
	      directions;
	\item \verb!#!\texttt{THAMES:Image{\us}Resolution} is the linear dimension of each cubic
	      voxel, in \unit{\micro\meter} units.
\end{itemize}

After the five-line header to this file, each line contains a single integer corresponding to the ID of
one of the microstructure phases (not the GEM phases).  The rows have a
\verb!z-y-x! nesting convention, by the first voxel is $(0,0,0)$, the $x$ coordinates
vary most quickly and the $z$ coordinates vary most slowly.

One way to generate such a file is to create a virtual 3D microstructure using VCCTL,
which produces a \verb!.img! file in much the same format as just shown.  However,
the VCCTL phase ID numbers (Table~\ref{tab:vcctlphases}) are different than the
ones required by THAMES.  Therefore, one must replace
the VCCTL ID numbers with the corresponding THAMES microstructure id numbers.
A C program called \verb!vcctl2thames! has been provided to assist with this, but the
program may need to be edited as needed to be consistent with the THAMES microstructure
definition file described later in this document.

\small{
	\begin{table}
		\caption{\label{tab:vcctlphases} VCCTL 9.5 phase identification numbers}
		\begin{tabular}{lSlSlS} \toprule
			Phase       & \multicolumn{1}{c}{ID} &
			Phase       & \multicolumn{1}{c}{ID} &
			Phase       & \multicolumn{1}{c}{ID}                                                             \\ \midrule
			Water       & 0                      & Gypsum                  & 7  & \ch{Ca(OH)2}          & 19 \\
			Void        & 55                     & Bassanite (Hemihydrate) & 8  & \ch{C-S-H}            & 20 \\
			Alite       & 1                      & Anhydrite (\ch{CaSO4})  & 9  & \ch{C3AH6}            & 21 \\
			Belite      & 2                      & Silica Fume             & 10 & Ettringite            & 22 \\
			\ch{C3A}    & 3                      & Inert                   & 11 & \ch{Fe(OH)3}          & 25 \\
			\ch{C4AF}   & 4                      & Silica glass            & 16 & Pozzolanic \ch{C-S-H} & 26 \\
			\ch{K2SO4}  & 5                      & Monosulfate             & 24 & Friedel salt          & 29 \\
			\ch{Na2SO4} & 6                      & Monocarboaluminate      & 34 & \ch{CaCl2}            & 28 \\
			\ch{CaCO3}  & 33                     & Str{\"{a}}tlingite      & 30 & Brucite               & 35 \\ \bottomrule
		\end{tabular}
	\end{table}
}

\subsection{\label{sec:thermodata} Thermodynamic data input}
THAMES 5.1 comes with the thermodynamic data files you will need to run most
of the simulations you want. This includes systems with temperatures ranging
from \qty{283}{\kelvin} to \qty{333}{\kelvin} and including most of the
important components of portland cement, portland limestone cement, and oil well
cement. It also includes most of the iron sulfides needed for simulations
involving pyrrhotite oxidation, as well as some common supplementary
cementitious materials (SCMs) including silica fume and proxies for the main
glassy and crystalline phases found in fly ashes. A full list of all the phases
included can be found in the appendix.

The thermodynamic data files can be found in any of the various test cases
provided with THAMES 5.1, located in subfolders of the \verb!tests! folder:
\begin{itemize}
	\item \verb!alite-298K-sealed-wc35!
	\item \verb!highvolume-silicafume!
	\item \verb!PC-Carbonate-200!
	\item \verb!PC-FlyAsh-200!
	\item \verb!portcem-298K-sat-wc45!
	\item \verb!portcem-298K-sealed-alk-wc45!
	\item \verb!portcem-298K-sealed-wc45!
	\item \verb!pyrragg-298K-sat!
	\item \verb!pyrrcem-296K-sat!
\end{itemize}
The first seven test cases in this list all have the exact same thermodynamic
data files, a collection of five files that all start with \verb!thames-!
The last two test cases have data files that contain chlorinated phases such
as Friedel's salt and Kuzel's salt, but they do not contain the glassy and
crystalline fly ash components.

If you need a system that cannot be modeled with these thermodynamic
data files, you will need to use the GEM-Selektor software specified in
Section~\ref{sec:requirements} and follow its documentation for creating
recipes and exporting data files.

\normalsize{ }
\subsubsection{\label{sec:thermodataformat} Meaning of GEMS3K the File Contents}
For normal usage of THAMES, one does not need to be concerned with most of the contents of
the GEMS3K files.  However, a few noteworthy items are as follows:
\begin{itemize}
	\item The \verb!-dat.lst! file contains that string that the GEMS3K library reads when
	      performing a calculation so that it knows which other files contain the thermodynamic
	      data it needs to run.
	\item The \verb!-dch.dat! file contains the names of the independent components
	      in the \verb!<ICNL>! field, the names of the dependent components in the
	      \verb!<DCNL>! field, and the names of the phases in the \verb!<PHNL>! field.
	      You will need these names when correlating microstructure phases with GEM phases.
	      This file also contains the parameters needed to calculate the molar volumes
	      and densities of the dependent components as a function of temperature and pressure.
	\item The \verb!-dbr.dat! file contains information about the calculated state of
	      the system calculated using GEM Selektor, such as the activity coefficients,
	      and phase abundances.  But these will change with subsequent calculations in
	      THAMES and so they are not of concern in the great majority of cases.
	\item The \verb!-ipm.dat! file contains information about the numerical parameters
	      used in the calculations, and are not of immediate concern for using THAMES.
\end{itemize}

%\fbox{
%\begin{minipage}{0.9\textwidth}
%    \textbf{WARNING}:  With the sole exception of the modification to the
%    \texttt{-sol-dat.lst} file already described, you must \textsc{not} modify
%    any of the other GEMS3K input files.  Doing so will likely cause the simulation
%    to crash with a GEM error.
%\end{minipage}
%}

\subsection{\label{sec:materialsystemfiles} Simulation Parameters}
\subsubsection{\label{sec:microphasedefs} Microstructure phase definition
	section}
THAMES defines a collection of phases that can potentially be present in its
3D microstructure. These are called ``(THAMES) microstructure phases,'' and
they can be defined as one or more of the phases recognized in the
thermodynamic data files from GEMS. For example, one microstructure phase
could be defined for each of the multiple phases in AFm family that
GEMS recognizes, or one microstructure phase could be defined as
the collection of all of those AFm phases. The user has the liberty to
define these microstructure phases, and it is not necessary to include
every phase recognized by GEMS within the collection of microstructure phases.
If you know that silica fume will never be present within your simulation,
it does not need to be defined in any microstructure phase that you create.
However, two non-negotiable microstructure phases \textbf{must} always be defined:
\textit{void space} and \textit{electrolyte} (the aqueous pore solution).

The other information that must be given is the final age of the system
to be simulated and the specific times that you want to output
microstructure images. These are both given in units of days.

The file for specifying these simulation parameters
in a JSON file, which you can name anything you want but which has a default
name of \verb!simparams.json!. The documentation for JSON syntax can be found at
\href{https://www.json.org}{https://www.json.org}.
We will examine the typical contents of this file for a portland
cement paste with some pyrrhotite included.

\subsubsection{\label{sec:envsection} Environment section.}
The first section is called \verb!environment! and contains
information about the curing conditions and any boundary conditions
on the electrolyte composition. An example is shown below, with
extra annotations that are not included in the actual JSON file

\scriptsize{
	\begin{lstlisting}
{
  "environment": {
    "temperature": 296.15,        # Curing temperature (K), required
    "reftemperature": 298.15,     # Reference temperature (K) for kinetics, required
    "saturated": 1,               # 1 = saturated pores, 0 = sealed pores,
    "electrolyte_conditions": [   # (optional)
      { "DCname": "Ca(CO3)@", "condition": "initial", "concentration": 1.0e-6 },
      { "DCname": "AlO2H@", "condition": "initial", "concentration": 1.0e-6 },
      { "DCname": "CaSiO3@", "condition": "initial", "concentration": 1.0e-6 },
      { "DCname": "Fe(CO3)@", "condition": "initial", "concentration": 1.0e-6 },
      { "DCname": "MgSiO3@", "condition": "initial", "concentration": 1.0e-6 },
      { "DCname": "KOH@", "condition": "initial", "concentration": 1.0e-6 },
      { "DCname": "Ca(SO4)@", "condition": "initial", "concentration": 1.0e-6 },
      { "DCname": "K+", "condition": "initial", "concentration": 2.0e-6 },
      { "DCname": "SO4-2", "condition": "initial", "concentration": 1.0e-6 },
      { "DCname": "Na+", "condition": "initial", "concentration": 0.672 },
      { "DCname": "ClO-", "condition": "initial", "concentration": 0.672 }
    ],
  },
  .
  .
  .
}
\end{lstlisting}
}

\normalsize{ }
The \verb!temperature!, \verb!reftemperature!, and \verb!saturated! lines
are required for every simulation. The \texttt{electrolyte{\us}conditions} section
is techincally optional, but it is a good idea to have it. The GEMS library
used by THAMES is more likely to properly converge if there is at least a little
bit of each independent component in the electrolyte at the beginning. As shown,
this can be done by including small amounts of different dependent components
(DCs) that contain the ICs.

The list of all available DCs can be found in the \verb!DCNL! section of the
thermodynamic DCH file (usually called \verb!thames-dch.dat!).
The required options for any electrolyte component are:
\begin{itemize}
	\item \verb!DCname!, must exacdtly match a name in the \verb!DCNL! section
	      of the \verb!-dch.dat! file.
	\item \verb!condition! can be ``initial'' if the component's concentration
	      is to be set initially but can vary as the simulation proceeds, or ``fixed''
	      if the component's concentration is to remain at the same value throughout.
	\item \verb!concentration! is the molal concentration for the component in
	      units of moles per kilogram of water.
\end{itemize}
Please be aware that the net electrical charge of the electrolyte must
be zero, and it is the user's reponsibility to ensure that the concentrations
are set in a way to balance the charge. If the charge is \textbf{not} balanced,
the simulation will throw an exception and exit.

\subsubsection{\label{sec:microphases} Microstructure phase definitions}
The next section is called \verb!microstructure! and defines each
of the phases (or collection of phases) that can potentially exist
within the 3D microstructure. The characteristics of a microstructure
phase definition are
\begin{itemize}
	\item its name,
	\item a unique numerical identification number (id) (\textbf{Note}
	      that each id appearing in the initial 3D microstructure image
	      file must be equal to the numerical id for one of the microstructure
	      phases defined here),
	\item whether or not it is considered to be part of the cementitious
	      material or is some other component in the mixture,
	\item its relation to phases defined in the thermodynamic data,
	\item its internal, or ``subvoxel'' porosity (assumed to be zero if
	      not given),
	\item its internal pore size distribution (optional),
	\item information about its kinetic behavior (more details to follow below),
	      and
	\item the color it should have in visualizations produced by THAMES
\end{itemize}

The section opens with its name (microstructure) and the number
of microstructure phases that are defined
\scriptsize{
\begin{lstlisting}
  "microstructure": {
    "numentries": 23,
    .
    .
    .
  },
\end{lstlisting}

\normalsize{ }
In this example there are going to be 23 microstructure phases
defined. Remember that every phase in the microstructure must be
defined here, but that not every phase defined here must appear
in the microstructure.
Immediately after the \verb!numentries!, each phase is
defined sequentially in a collection named \verb!phases!.

\paragraph{The void phase.} The first defined phase \textit{must} be ``Void'' which
is just empty porosity within the microstructure. This phase is unique
because it is the only one that does not have any associated thermodynamic
phase. The meaning of each entry is
\scriptsize{
	\begin{lstlisting}
    "phases": [
      {
        "thamesname": "Void",
        "id": 0,
        "cement_component": 0,
        "display_data": { "red": 0.0, "green": 0.0, "blue": 0.0, "gray": 0.0 }
       },
       .
       .
       .
    ],
  },
\end{lstlisting}
}

\normalsize{ }
The meaning of each entry is
\begin{itemize}
	\item \verb!thamesname! is any string in quotes to designate the phase.
	      The void phase \textbf{must} be named \verb!Void! and the pore
	      solution phase \textbf{must} be named \verb!Electrolyte!, but
	      you are free to choose the names of all other microstructure phases.
	\item  \verb!id! is the unique numerical identification number, which
	      must be a non-negative integer. The void phase \textbf{must} have
	      \verb!id! of 0 and the pore solution \textbf{must} have
	      \verb!id! of 1, but you are free to choose the \verb!id!
	      of all other microstructure phases. Just remember that the phase
	      identificatino numbers in the 3D microstructure image file must
	      correspond to the correct microstructure phase \verb!id! here.
	\item \texttt{cement{\us}component} is 1 for any phase that is solid and
	      that belongs to the cementitious material, such as cement clinker phases, gypsum,
	      alkali sulfates, free lime, \textit{etc}, and is 0 for other
	      phases such as void, electrolyte, silica fume, fly ash phases, and
	      any solid hydration products such as \ch{C-S-H} gel. This distinction
	      is needed for internal calculations of water/cement ratio and
	      water/solids ratio.
	\item \texttt{display{\us}data} is a set of red/green/blue/gray values
	      that specify the color of the phase in THAMES visualizations. Many
	      microstructure phases, though not all, have default colors assigned
	      already according to predefined \verb!thamesname! names, so this entry
	      is only needed if you are defining a new
	      microstructure phase, changing the name of an existing phase,
	      or if you wish to override the default color.
\end{itemize}

\paragraph{The electrolyte phase.} The next phase after the void phase
\textbf{must} be the pore solution, which has the name \verb!Electrolyte!:
\scriptsize{
	\begin{lstlisting}
  {
    "thamesname": "Electrolyte",
    "id": 1,
    "cement_component": 0,
    "gemphase_data": [
      {
        "gemphasename": "aq_gen",
        "gemdc": [
          { "gemdcname": "Al(SO4)+", "gemdcporosity": 1.0 },
          { "gemdcname": "Al(SO4)2-", "gemdcporosity": 1.0 },
          { "gemdcname": "Al+3", "gemdcporosity": 1.0 },
          { "gemdcname": "AlO+", "gemdcporosity": 1.0 },
          { "gemdcname": "AlO2-", "gemdcporosity": 1.0 },
          { "gemdcname": "AlO2H@", "gemdcporosity": 1.0 },
          .
          .
          .
        ],
      }
    ]
  }
\end{lstlisting}
}

\normalsize{ }
This phase definition starts just like the void phase, with a
\verb!thamesname!, a unique \verb!id!, and the designation
that it is \textbf{not} to be considered a cement component.
After these three, there is a section called \texttt{gemphase{\us}data}
that tells THAMES how
to relate this microstructure phase to one or more phases in
the thermodynamic system definition (\verb!-dch.dat! file).
\begin{itemize}
	\item \verb!gemphasename! is the name of the corresponding
	      GEMS system phase found in the \verb!PHNL! section of the
	      \verb!-dch.dat! file. It must match exactly.
	\item \verb!gemdc! is a list of all the GEMS dependenct
	      components (DCs) that are assigned to this GEM phase.
	      Each entry has two characteristics:
	      \begin{itemize}
		      \item \verb!gemdcname! is the name of the DC, which must
		            actually belong to the GEM phase and which must also
		            exactly match the name in the \verb!DCNL! section of
		            the \verb!-dch.dat! file. Note that these correspondences
		            of GEM phases and GEM DCs are already set up in the
		            \verb!-dch.dat! file, but they are listed here primarily
		            for solid solutions where the internal pore space of
		            a phase may depend on its composition (see below).
		      \item \verb!gemdcporosity! is the volume fraction of
		            non-solid space associated with this DC. It must be a
		            real number on the interval $\left[0,1\right]$.
	      \end{itemize}
\end{itemize}
Note that the electrolyte phase has many DCs, one for each
possible dissolved component; only a small fraction of them are
shown here for brevity.

The names of all
the DC components for a GEM phase must be placed within \verb!gemdc! field,
and there must be one of these fields for each GEM phase that you wish to link
to the microstructure phase.  In addition, the \verb!gemphasename! must exactly match
one of the \verb!<PHNL>! entries in the \verb!-dch.dat! file, and
\textit{all} of the DCs associated with that phase must immediately follow that
phase as shown in this example.  Determine which GEM DCs are associated
with a given phase by consulting the \verb!<nDCinPH>! field of the \verb!-dch.dat! file.

To illustrate, Figure~\ref{fig:dchfile} shows a portion of a real \verb!-dch.dat! file.
In the figure, there are \textit{nine} phases listed in the \verb!<PHNL>! field and
26 DCs listed in the \verb!<DCNL>! field.  The numbers in the \verb!<nDCinPH>! field
relate the two.  Notice that there are \textit{nine} numbers in \verb!<nDCinPH>!, one
for each phase. The ordering of the numbers is the same as the ordering of phase
names in the \verb!<PHNL>! field.  So there are 13 DCs in the first phase, which
is \texttt{aq{\us}gen}.  In addition, those 13 DCs are the \textit{first} 13 DCs in the
\verb!<DCNL>! list.  Moving on, we see that there are three DCs in the next phase, which
is \texttt{gas{\us}gen}, and those three DCs are the next three after the first 13 in the
\verb!<DCNL>! list, namely \verb!'H2'!, \verb!'O2'!, and \verb!'H2O'! (the last being
the water vapor molecule, not liquid water which is \verb!H2O@!\footnote{In GEMS, the
	\makeatletter\texttt{`@'}\makeatother\ symbol is reserved to designate
	a zero-charge liquid solvent or
	a zero-charge dissolved component.  Gas molecules do not have this symbol even though
	they also have zero charge.}

\begin{figure}
	\centering\includegraphics[width=0.7\textwidth]{../Figures/dchfile.png}
	\caption{\label{fig:dchfile} Portion of a DCH input file showing the definitions
		of GEM phases and DCs.}
\end{figure}

\paragraph{Solid microstructure phases.} Solid microstructure
phases may be ``simple'' (one phase, one DC, and possibly
some impurities), ``solid solution'' (one phase, multiple DCs
that may be present in any mole fraction to give a phase
with nonstoichiometric composition, such as \ch{C-S-H} gel),
or ``composite'' (multiple thermodynamic phases combined for
simplicity or computational efficiency).

An example of a ``simple'' microstructure phase is alite,
which is usually a monoclinic form of tricalcium silicate
(\ch{Ca3SiO5}) with some impurities to stabilize the monoclinic
polymorph.
\scriptsize{
	\begin{lstlisting}
  {
    "thamesname": "Alite",
    "id": 2,
    "cement_component": 1,
    "interface_data": { }
    "gemphase_data": [
      {
        "gemphasename": "Alite",
        "gemdc": [{ "gemdcname": "C3S" }]
      }
    ],
    "impurity_data": {
      "k2ocoeff": 0.00087,
      "na2ocoeff": 0.0,
      "mgocoeff": 0.00861,
      "so3coeff": 0.007942
    },
    "kinetic_data": {
      .
      .
      .
    }
  },
 \end{lstlisting}
}

\normalsize{ }
This entry indicates three more features of a solid phase, namely
data about its interface with other phases, its impurity content,
and its kinetic behavior:
\begin{itemize}
	\item \texttt{interface{\us}data} in the above entry for alite
	      is empty, but can have information about the affinity the
	      phase has for growing on one or more other substrate phases, using
	      the name of the substrate phase and the thermodynamic contact
	      angle between the two:
	      \begin{itemize}
		      \item \verb!affinityphase! must be the \verb!thamesname! of
		            another \textit{microstructure} phase defined in this file
		      \item \verb!contactangle! is the equivalent thermodynamic
		            contact angle for the two phases in units of degrees,
		            and must be a real number on the interval $\left[0,180\right]$,
		            where \qty{0}{\degree} is perfect wetting and
		            \qty{180}{\degree} is complete nonwetting.
	      \end{itemize}
	\item \texttt{impurity{\us}data} contains data on the amounts
	      of four different possible impurities, namely
	      potassium, sodium, magnesium, and sulfur, in terms
	      of the \textit{mole fractions} of their equivalent oxides,
	      namely \ch{K2O}, \ch{Na2O}, \ch{MgO}, and \ch{SO3}.
	      Currently these are the only four impurities that
	      are recognized by THAMES.
	\item \texttt{kinetic{\us}data} is optional and will be described
	      in detail in a later section. For now, the presence of this
	      section indicates that the rates of reaction of this phase
	      will be controlled by a kinetic rate equation. If this
	      section is not provided, the amounts of this phase at any
	      time will be determined solely by thermodynamic equilibrium
	      as calculated by GEMS.
\end{itemize}

An example of a more complicated composite microstructure phase could
be a collection of carbonated AFm phases, which would include all the
monocarbonates, hemicarbonates, and tricarboaluminate phases:
\scriptsize{
	\begin{lstlisting}
      {
        "thamesname": "AFmc",
        "id": 14,
        "cement_component": 0,
        "interface_data": {
          "affinity": [
            { "affinityphase": "Alite", "contactanglevalue": 180 },
          ]
        },
        "gemphase_data": [
          {
            "gemphasename": "C4Ac0.5H105",
            "gemdc": [{ "gemdcname": "hemicarb10.5" }]
          },
          {
            "gemphasename": "C4Ac0.5H12",
            "gemdc": [{ "gemdcname": "hemicarb" }]
          },
          {
            "gemphasename": "C4Ac0.5H9",
            "gemdc": [{ "gemdcname": "hemicarb9" }]
          },
          {
            "gemphasename": "C4AcH11",
            "gemdc": [{ "gemdcname": "monocarb" }]
          },
          {
            "gemphasename": "C4AcH9",
            "gemdc": [{ "gemdcname": "monocarb9" }]
          }
        ]
      },
 \end{lstlisting}
}

\normalsize{ }
The absence of impurities or kinetic data in this example indicate
that the phase is to be considered pure and that its amount in the
microstructure at any time will be determined solely by thermodynamic
equilibrium with the electrolyte. Also note that each of the DCs
associated with this phase default to having zero porosity, so there
will be no subvoxel or interhydrate porosity associated with this
microstructure phase.

Finally, a solid solution with internal porosity can be defined, such
as \ch{C-S-H} gel:
\scriptsize{
	\begin{lstlisting}
      {
        "thamesname": "CSHQ",
        "id": 11,
        "cement_component": 0,
        "interface_data": {
          "affinity": [
            { "affinityphase": "Alite", "contactanglevalue": 30 },
            { "affinityphase": "AFmc", "contactanglevalue": 90 },
          ]
        },
        "gemphase_data": [
          {
            "gemphasename": "CSHQ",
            "gemdc": [
              { "gemdcname": "CSHQ-JenD", "gemdcporosity": 0.4935 },
              { "gemdcname": "CSHQ-JenH", "gemdcporosity": 0.4935 },
              { "gemdcname": "CSHQ-TobD", "gemdcporosity": 0.2004 },
              { "gemdcname": "CSHQ-TobH", "gemdcporosity": 0.2004 },
              { "gemdcname": "KSiOH", "gemdcporosity": 0.1825 },
              { "gemdcname": "NaSiOH", "gemdcporosity": 0.1825 }
            ]
          }
        ],
        "poresize_distribution": [
          { "diameter": 0.23228, "volumefraction": 0.01684781722375019 },
          { "diameter": 0.49089, "volumefraction": 0.02276126023357459 },
          .
          .
          .
        ],
        "Rd": [
          { "Rdelement": "K", "Rdvalue": 0.42 },
          { "Rdelement": "Na", "Rdvalue": 0.42 }
        ]
      },
\end{lstlisting}
}

\normalsize{ }
Two more sections can be found in this example. First, the pore size
distribution can be defined as a normalized probability density function,
with the differential volume fraction (\verb!volumefraction!, dimensionless)
given for each effective pore size dimension (\verb!diameter!) given in
nanometer units. The second new section is the impurity partition section,
with the name \verb!Rd!, that describes the proportion of impurities that will
be drawn from the surrounding electrolyte if available:
\begin{itemize}
	\item \verb!Rdelement! is the name of an independent component (IC)
	      found in the \verb!ICNL! section of the \verb!-dch.dat! file.
	\item \verb!Rdvalue! is the partition coefficient defined for
	      alkalis by Hong and Glasser (\textit{Cem Concr Res}, 29 (1999), 1893--1903;
	      \textit{Cem Concr Res}, 32 (2002), 1101-1111.) which is defined
	      \begin{equation}
		      R_\text{d} = \frac{c_s w}{c_d s} \qquad (\text{units of}\
		      \unit{\liter\per\kilo\gram})
	      \end{equation}
	      where $c_s$ is the concentration of the impurity in the solid phase
	      (units \unit{\mole\per\liter}), $c_d$ is the concentration of the impurity
	      in the electrolyte (units \unit{\mole\per\liter}), and $w/s$ is the
	      water-solids ratio in units of \unit{\liter\per\kilo\gram}
\end{itemize}

\subsubsection{\label{sec:timeparams} Calculation and output times}
The final section of this file is called \texttt{time{\us}parameters} and defines
both the final age of the material to simulate and the simulation times at
which THAMES should write data for visualization and for assessing the
pore size distribution and saturation state.

\scriptsize{
	\begin{lstlisting}
  "time_parameters": {
    "finaltime": 56.0,
    "outtimes": [0.01, 0.1,0.5,1.0,3.0,7.0,14.0,28.0]
  }
\end{lstlisting}
}

Or, alternatively,

\scriptsize{
	\begin{lstlisting}
  "time_parameters": {
    "finaltime": 56.0,
    "outfreq": 4.0,
  }
\end{lstlisting}
}

\normalsize{ }
The meaning of the parameters are as follows:
\begin{itemize}
	\item \verb!finaltime! is the final time, in days, to attempt simulating
	      the microstructure evolution. For various reasons the simulation may not
	      complete the full requested age---for example, the microstructure could
	      be completely depleted of electrolyte and so not be able to perform any
	      other reactions---but the simulation will definitely stop
	      after this time is reached.
	\item \verb!outtimes! is a list of times, in days, at which THAMES should
	      output the microstructure state, which includes
	      \begin{itemize}
		      \item the full 3D microstructure image (an \verb!.img! file),
		      \item a 2D image of a slice of the microstructure (both \verb!.ppm!
		            and \verb!.png! files), and
		      \item An entry in the \verb!.xyz! file used for 3D
		            visualizations.\footnote{The 3D visualization option is
			            discussed more in Section~\ref{sec:3Drendering}.}
		      \item the current pore size distribution and saturation state.
	      \end{itemize}
	\item \verb!outfreq! is an alternative to \verb!outtimes! offered for
	      convenience. It allows one to specify a time interval, in days, at the end
	      of which the output images and pore size distributions will be output.
	      This enables one to avoid entering many equally spaced times for output.
	      Note: If both \verb!outfreq! and \verb!outtimes! are present, the
	      \verb!outtimes! entry will be ignored.
\end{itemize}

\section{\label{sec:outputguide} THAMES Output Files}
This section provides guidance on the information contained in the various output files
produced by THAMES v5.1. For simplicity, the guide will be structured entirely
in the form of an example simulation based on the test case named
\verb!portcem-296K-sealed-alk-wc45! that is provided.

\section{\label{sec:assumption} Assumptions}
This guide assumes
\begin{enumerate}
	\item that you have been through the previous sections of this document,
	\item that your THAMES 5.1 distribution located in a directory
	      with the absolute path specified in this guide by \verb!$THAMES!. For example
	      if your distribution is located at the absolute path
	      \verb!/home/odehwah/models/THAMES!, then \verb!$THAMES! is equivalent to that
	      full path in this guide; and
	\item that you have already built \verb!thames! and that it is in
	      your \verb!$PATH! environment variable so that the operating system can discover it
	      when you run it from the command line. If you are unsure whether this has been
	      done, you can run the command \verb!which thames! from the command line. If that
	      command returns something like \verb!thames not found!, then you must first
	      augment your \verb!$PATH! environment variable. If you are unsure how to
	      change the environment variable a tutorial for Linux can be found
	      here: \url{https://www.digitalocean.com/community/tutorials/how-to-view-and-update-the-linux-path-environment-variable}
\end{enumerate}

\section{\label{sec:runtestcase} Run the Test Case}
A good idea is to keep the working directory as clean as possible so that it
doesn't get accidentally corrupted or more confusing. So it is a good practice
to copy the test case to another location before using it there. For this
example, create a temporary directory and then copy the test case there:
\scriptsize{
	\begin{lstlisting}
  mkdir $HOME/scratch
  cp -R $THAMES/test/portcem-296K-sealed-alk-wc45 $HOME/scratch
\end{lstlisting}
}

\normalsize{ }
Next, go to that new directory:
\scriptsize{
	\begin{lstlisting}
  cd $HOME/scratch/portcem-296K-sealed-alk-wc45
\end{lstlisting}
}

\normalsize{ }
Before running the simulation, open the \verb!input.in! file. The file contents
are:
\scriptsize{
	\begin{lstlisting}
2
thames-dat.lst
simparams.json
ccr140-w45-thames.img
cem140-sealed-alk-01
\end{lstlisting}
}

\normalsize{ }
Each line of this file is now explained:
\begin{itemize}
	\item \textbf{Line 1}: Determines the type of simulation to be run:
	      \begin{enumerate}
		      \item A value of 1 exits the simulation with no further action
		      \item A value of 2 will run a normal hydration simulation
	      \end{enumerate}
	\item \textbf{Line 2}: The root name of the thermodynamic data. This file is
	      read and automatically prompts the input of the other thermodynamic data
	      files.
	\item \textbf{Line 3}: The JSON formatted input file specifying the simulation
	      environment, the microstructure phase definitions, and the simulation time
	      and output frequency.
	\item \textbf{Line 4}: The name of the file containing the initial 3D microstructure
	\item \textbf{Line 5}: The root name for all output files produced by THAMES.
	      This can be anything you want, although it is recommended not to use periods
	      within it. Feel free to shorten it to \verb!cem140! or some other string if
	      you wish.
\end{itemize}
Close the file, making sure to save it first if you made changes to the fifth
line, and then run the THAMES simulation with these options:
\scriptsize{
	\begin{lstlisting}
  thames --xyz --outfolder MyResult < input.in >& output.out &
\end{lstlisting}
}

\normalsize{ }
The simulation should create an \verb!output.txt! file and a folder named
\verb!MyResults!, as shown in Figure~\ref{fig:one}, while it continues
to run for about five minutes to ten minutes.

\begin{figure}
	\centering\includegraphics[width=0.7\textwidth]{../Figures/ListingAfterRun.png}
	\caption{\label{fig:one} Contents of the test directory running a simulation.}
\end{figure}

A message similar to the following
appear in the terminal when the simulation finishes:
\scriptsize{
	\begin{lstlisting}
  [1]  + 56549 exit 1     ~/Software/THAMES/bin/thames --xyz --outfolder MyResult < input.in >&
\end{lstlisting}
}

\normalsize{ }
The \verb!MyResults! folder
holds all of the output files created by the simulation, while the
\verb!output.txt! file holds all the standard output that would otherwise
be written to the screen. Under ordinary circumstances there is no need to
examine the \verb!output.txt! file, but it can be useful to have it if
the program crashes for some reason or otherwise behaves oddly.

\section{\label{sec:outputfiles} Output File Content}
When the simulation exits, navigate to the \verb!MyResults! folder and
run the \verb!ls! command to see the directory contents, which should
appear like that shown in Fig.~\ref{fig:two}.

\begin{figure}
	\centering\includegraphics[width=0.7\textwidth]{../Figures/OutputFileTypes.png}
	\caption{\label{fig:two} Contents of the \texttt{MyResult} folder at the
		conclusion of THAMES simulation. Color blocks superimposed on the files
		indicate spreadsheet-like CSV data (green), pore size distribution data
		(yellow), microstructure data and images for viewing (pink), a 3D
		\texttt{xyz}
		format file for viewing in Ovito or similar application (orange, optional),
		and original input data files (blue) copied from the parent directory for
		record-keeping.}
\end{figure}

\subsection{Input data files}
The files highlighted in blue in Fig.~\ref{fig:two} are input files
copied directly from the parent directory and are present only as
metadata to recall the conditions under which the simulation was
performed. The contents of most of these files are described in detail
in the companion document on preparing input files and will therefore
not be revisited here.

\subsection{CSV files}
Eleven comma separated value CSV files are written by THAMES to
recored time-dependent simulation data. These are highlighted in green in
Fig.~\ref{fig:two} and are described in the following paragraphs:

\paragraph{{\us}CSH.csv} Records the composition of the calcium silicate
hydrate (\ch{C-S-H}) phase as a function of time. The first column is the time (in hours).
The rest of the columns up until the final column are the concentrations
of each element in the \ch{C-S-H} phase. The final column displays the
molar Ca/Si ratio of the \ch{C-S-H}.

\paragraph{{\us}CSratio{\us}solid.csv} Provides the overall molar Ca/Si ratio of all
solid components in the microstructure. Time (in hours) is in the first column
and the only other column is the averaged Ca/Si ratio. This file may be
removed in future versions of THAMES.

\paragraph{{\us}dcmoles.csv} Catalogs the number of moles (per \qty{100}{\gram}
of initial solid) of each dependent component (DC)\footnote{Note that DCs are
	not identical with phases. Some phases, such as gypsum, have only one
	DC, which for gypsum is named \texttt{Gp}, while phases such as \ch{C-S-H}
	are solid solutions of multiple DCs.}
in the system, again with
simulation time (in hours) occupying the first column.

\paragraph{{\us}DCVolumes.csv} Provides the volume (in \unit{\meter\cubed} per
\qty{100}{\gram} of initial solid) of each dependent component (DC) in the
system, again with simulation time (in hours) occupying the first column.

\paragraph{{\us}Enthalpy.csv} Displays in the second column the predicted cumulative heat release
by the system (in \unit{\joule} per \qty{100}{\gram} of initial solid) at each
corresponding time (in hours) in the first column.

\paragraph{{\us}icmoles.csv} Gives the calculated moles, per \qty{100}{\gram}
of initial solid, of each independent
component (IC, essentially chemical elements and electric charge) in the system
for each corresponding time (in hours) in the first column. This file may be
removed in future versions of THAMES because it duplicates information that
can be easily obtained from the \texttt{{\us}dcmoles.csv} file.

\paragraph{{\us}Microstructure.csv} Records the \textit{volume fraction}
(dimensionless) of each defined microstructure phase in the system
at the corresponding time (in hours) displayed in the first column.
The very last column gives the predicted chemical shrinkage value at the
corresponding time, in units of \unit{\meter\cubed} per \qty{100}{\gram}
of initial solid material.

\paragraph{{\us}SI.csv} Contains the saturation index of the pore solution
(dimensionless) with respect to each defined microstructure phase
at each corresponding time (in hours) in the first column. Here the saturation
index, $\Omega$, is defined as the ratio of the activity product to the equilibrium
constant, so for any given solid phase $p$, $\Omega_p = 1$ if the solution is at
equilibrium with that solid phase, $\Omega_p < 1$ if the solution is
undersaturated so that the phase $p$ is thermodynamically driven to dissolve,
and $\Omega_p >$ if the solution is
supersaturated so that the phase $p$ is thermodynamically driven to precipitate.

\paragraph{{\us}Solution.csv} Displays the composition of the electyrolyte
(\textit{i.e.}, pore solution) at each calculated time (in hours) given in
the first column. The composition consists of the concentration of each
dependent component dissolved in the water, given in units of
\unit{\mole\per\kilo\gram} water.

\paragraph{{\us}SurfaceAreas.csv} Displays the solid-water surface
area of each defined microstructure phase at each calculation time (in hours)
displayed in the first column. The surface areas in this file are given
in units of \unit{\meter\squared} per \qty{100}{\gram} of initial solid
material. In particular, they are not equivalent to the specific surface
area of each phase individually, but come closer to reflecting each phase's
contribution to the specific surface area.

\subsection{Pore size distribution data files}
One of these files is created for each output time the user specifies in the
input JSON file, which is described in more detail in the companion guide
for preparing input files. Each file has a name that contains the
simulation temperature in kelving units and the time in year-day-hour-minute format.
For example, a file with the name
\scriptsize{
	\begin{lstlisting}
myfile_PoreSizeDistribution.001y.056d.14h.30m.296K.csv
\end{lstlisting}
}
\normalsize{ }
corresponds to a temperature of \qty{296}{\kelvin} and
a simulation time of 1 year, 56 days, 14 hours, and 30 minutes.
Times are always rounded to the nearest minute.

The contents of a pore size distribution file are structured
as shown below:
\scriptsize{
	\begin{lstlisting}
Time = 672 h,,
Voxel-scale pore volume fraction (>= 1000 nm) = 0.049738,,
Saturated voxel-scale pore volume fraction = 0,,
Empty voxel-scale void volume fraction = 0.049738,,
Subvoxel volume fraction (< 1000 nm) = 0.173083,,
Saturated subvoxel volume fraction (< 1000 nm) = 0.13513,,
Empty subvoxel-scale void volume fraction = 0.0379536,,
Total pore volume fraction = 0.222821,,
Total void volume fraction = 0.049738,,
Pore size saturation data:,,
Masterporevolume size = 56,,
Diameter (nm),Volume Fraction,Fraction Saturated
1.99526,0.0426508,1
2.23872,0.00530618,1
2.51189,0.00477843,1
2.81838,0.00436585,1
    .
    .
    .
63.0957,0.00507889,1
70.7946,0.00853574,0.345385
79.4328,0.00906405,0
89.1251,0.00953302,0
100,0.0115934,0
223.872,0.0021721,0
501.187,3.43993e-06,0
>=1000,0.049738,0
\end{lstlisting}
}

\normalsize{ }
The first eleven lines provide overall data on the volume fraction of different
ranges of saturated pores and empty pores. Following this is the more detailed
analysis of pore volume fractions. The first column is the effective
pore diameter in  \unit{\nano\meter} units, the second column is the
volume fraction (dimensionless) of the microstructure occupied by pores of that size,
and the third column is the fraction of that pore volume occupied by the
pore solution. The data are provided up to a maximum pore diameter of
\qty{99}{\nano\meter}, and all larger pores are lumped into a final category
on the final row.

\subsection{Microstructure state files}
For each specified output time in the JSON input file, two or three files
will be created:
\begin{itemize}
	\item \verb!.img! files contain the full 3D microstructure image data
	      in exactly the same format as the input microstructure image, which
	      is discussed in detail in the companion guide for preparing input
	      files.
	\item \verb!.ppm! is a 2D $xy$-slice through the microstructure at the
	      midpoint of the $z$ dimension. The format is a color portable pixel
	      map (PPM) that can be displayed by a range of image rendering applications
	      such as Gimp (any platform, most flexible), \verb!eog! or \verb!feh! (linux),
	      Preview (Mac OS), Honeyview or i{\us}view (Windows).
	\item \verb!.png! is a PNG version of the corresponding \verb!.ppm! file
	      that THAMES will generate automatically only if ImageMagick 7 or later
	      \url{https://imagemagick.org/} is
	      installed on the computer and if the \verb!magick! command is in the
	      path.
\end{itemize}

The \verb!.ppm! and \verb!.png! files are created with some sense of
depth by making the pore space slightly transparent. These two files
also render each microstructure phase with a unique color. The
\verb!rgb! triplet of values for the color of each phase is given
in the \texttt{{\us}Colors.csv} file for easier reference.

\subsection{\label{sec:3Drendering} 3D rendering file}
This file tends to be quite large and so it is only created if the user gives
the \verb!--xyz! option on the command line when starting a THAMES simulation.
It contains a full 3D data structure for each output time specified in the JSON
input file, and so can be rendered into individual states or produced into
an animation of the microstructure's time evolution. The file is written
in the industry-standard \verb!.xyz! format that can be read by multiple
3D rendering software applications. Among these,
Ovito (\url{https://ovito.org})\footnote{Mention of commercial products is
	for information only, and does not imply recommendation or endorsement by Texas A\&M University.}
is a freely available application that runs on any platform and is relatively easy to use.

\end{document}
